\section{Zusammenfassung und Entscheidungsvorlage}

Wenn eine moderne Produktionsanlage stehen bleibt, führt die herkömmliche Fehlersuche zu zeitraubenden Suchen von Maschinendaten und komplexen Analysen von Signalzuständen. 

Das in dieser Arbeit vorgestellte Konzept löst dieses Problem durch einen intelligenten Fehlersupport: Die Steuerung speichert im funktionsfähigem Zustand Aktionen und Messwerte in einem Ringspeicher.
Im Fehlerfall  werden die Daten automatisch von einer Künstlichen Intelligenz analysiert. Die KI zeigt dem Techniker Ursache und Behebung des Fehlers direkt auf dem HMI-Bedienbildschirm an.

Dadurch verkürzt sich die Stillstandszeit um durchschnittlich zehn Minuten pro Fehlerfall. Bei einmaligen Entwicklungskosten von rund 4.500 Euro spart das System über 26.000 Euro pro Jahr ein und amortisiert sich nach rund zwei Monaten.

\subsection{Entscheidungsfrage}
Soll das Konzept für den intelligenten Fehlersupport genehmigt, im Rahmen eines Studentenprojekts an einer realen Produktionslinie getestet und bei nachgewiesenem Erfolg als neuer Standard für alle weiteren Anlagen bei RAFI eingeführt werden?

\subsection{Konkrete Entscheidungsvorlage für das Management}
Dem Management wird empfohlen, das Budget von 4.500 Euro freizugeben, um das System innerhalb von acht Wochen an einer Testanlage in Betrieb zu nehmen. Während dieser Phase soll überprüft werden, ob die Reparaturzeiten um mindestens 25 Prozent sinken und die Praxistauglichkeit gegeben ist. Bei Erfolg wird die Lösung fest in den RAFI-Standard übernommen. Das Management beschließt damit, die Digitalisierung der Fertigung voranzutreiben, ungeplante Stillstände zu reduzieren und die Wirtschaftlichkeit nachhaltig zu steigern.